\section*{Chapitre \ref{ch:hf:a-venue-playbook} --- Un playbook de plateforme}

\begin{solution}{exo:hf:a-venue-playbook:1}
$0.125\times1\,200+5=155$ ordres. Les simulations dont le rythme d'envoi est de huit secondes ou moins ont accepté 151 à 154 ordres~: le budget est entièrement consommé, quoi qu'on envoie.
\end{solution}

\begin{solution}{exo:hf:a-venue-playbook:2}
Deux ordres toutes les $2/0.125=16$ secondes. La cotation sur la plateforme d'origine n'a aucun rythme (un écart nul), ce que refuse la validation du \omterm{def:hf:a-venue-playbook:profile}{profil de plateforme}.
\end{solution}

\begin{solution}{exo:hf:a-venue-playbook:3}
L'ordre de tête reçoit $\min(0.4\times1\,000,\,200)=200$. Les 800 autres contrats se partagent entre les 4\,800 restants au prorata de la taille~: le
teneur de marché reçoit $\lfloor 500\times800/4\,800\rfloor=\lfloor83.3\rfloor=83$, et les restes d'arrondi vont dans l'ordre FIFO.
\end{solution}

\begin{solution}{exo:hf:a-venue-playbook:4}
Chaque message de chaque agent est retardé des mêmes 350 microsecondes, de sorte que l'ordre de l'arbitragiste arrive toujours avant l'annulation du teneur de marché, et le
flux du tape, dont les événements sont espacés de dizaines de millisecondes, rencontre les ordres du teneur de marché dans le même ordre que sans le retard. Un retard asymétrique
qui épargne les annulations, ou un teneur de marché qui réagit à un signal que l'arbitragiste n'a pas, modifierait les courses entre ordres.
\end{solution}

\begin{solution}{exo:hf:a-venue-playbook:5}
Le retard protège les ordres au repos de l'arbitragiste, mais s'applique aux propres nouveaux ordres du teneur de marché~: 52~\% de son volume y a été consommé, et non fourni,
parce qu'un ordre coté pour le carnet qu'il voyait arrivait cinq secondes plus tard dans un carnet qui avait bougé et l'avait croisé. Ces prises se font à des cours périmés, dans
le sens du mouvement qu'avait déjà pris le cours.
\end{solution}

\begin{solution}{exo:hf:a-venue-playbook:6}
Dans une adjudication à prix uniforme, le meilleur cours acheteur du teneur de marché ne s'exécute que si le prix de règlement se situe à ce cours ou en dessous, ce qui arrive lorsque les ventes
de l'adjudication l'emportent sur ses achats~: ses exécutions surviennent lorsque le prix va le moins dans son sens. Il ne se fait jamais capturer, car son annulation sur saut arrive dans le même intervalle
que l'ordre de l'arbitragiste et est traitée avant le règlement. Le remède tient au calcul du prix, et non aux paramètres~: coter autour d'une fair value
anticipant le prix de règlement de l'adjudication, et dimensionner selon le déséquilibre.
\end{solution}

\begin{solution}{exo:hf:a-venue-playbook:7}
Un écart de huit secondes~: 10\,100 unités contre 8\,333 à seize secondes, avec 28 ordres rejetés en vingt minutes. La règle des seize secondes laisse
des jetons inutilisés (118 sur environ 155 acceptés)~; le rythme de huit secondes les utilise tous, en comptant sur la rafale pour absorber les moments où les deux côtés
ont besoin d'une recotation au même instant, et il le paie de quelques rejets.
\end{solution}

\begin{solution}{exo:hf:a-venue-playbook:8}
La validation vérifie le profil selon ses propres règles, et le simulateur vérifie la stratégie par rapport à son modèle de la plateforme~: aucun des deux
ne vérifie le profil par rapport à la plateforme. Les premiers jours doivent se dérouler dans l'environnement de test de la plateforme, puis en petite taille sous les limites de risque, en comparant
les exécutions, les rejets et les frais au profil, avant de passer à la pleine taille.
\end{solution}

\begin{solution}{pb:hf:a-venue-playbook:1}
\begin{enumerate}
\item Voir la \cref{def:hf:a-venue-playbook:profile}~; la moitié plateforme (ce que fait la plateforme) et la moitié cotation (ce que fait le desk sur celle-ci).
\item Appariement (au \cref{ch:hf:futures-market-making}), tick et lot, frais et unité (au \cref{ch:hf:rebates-inverted-venues-and-tiers}), flux
(au \cref{ch:hf:the-business-of-market-making}), types d'ordres (au \cref{ch:hf:the-quoting-engine}), retards (au \cref{ch:hf:latency-arbitrage-and-its-defence}),
budgets (au \cref{ch:hf:crypto-high-frequency-trading-on-centralised-venues}), séance et particularités.
\item Flux market-by-order avec de nouvelles références au remplacement~; un retard symétrique de 350 microsecondes et un retard asymétrique désapprouvé~; ordre de tête puis au
prorata sur certains contrats à terme~; ratios d'efficacité des messages~; outils de risque fondés sur l'exécution pour les options~; 100 ordres par 10 secondes et 200\,000 par jour sur une
plateforme crypto~; un retard en cours de jeu de une à douze secondes sur la prise de paris.
\item Pour que la liste des plateformes soit complète et qu'une stratégie supposant un carnet d'ordres soit refusée là où il n'y en a pas.
\item Dimensionner un multiple du lot~; un ordre post-only uniquement là où il est proposé~; des ordres post-only là où un retard les épargne~; recotations rythmées dans le budget d'ordres~;
annulations gratuites sous un budget.
\item Le budget crypto (sans rythme) et le lot FX (100 unités contre un million).
\item L'annulation sur saut est la protection du teneur de marché~; si elle attendait des jetons, un dépassement de budget se transformerait en capture.
\item Par les références contenues dans ses accusés de réception, qui arrivent avant le flux de données de marché~; si le flux arrivait en premier, le teneur de marché prendrait son propre ordre pour
celui d'un autre et le poursuivrait.
\item Le flux du tape du livre~7 rejoué sous forme d'ordres réels~; un saut du prix efficient toutes les 20 secondes en moyenne~; un arbitragiste qui en est informé le premier
et prend le meilleur niveau périmé~; le teneur de marché l'apprend 20 microsecondes plus tard et annule son côté périmé.
\item Contrats à terme transparents, inversés, à \omterm{def:hf:latency-arbitrage-and-its-defence:bump}{ralentisseur} et FIFO~: 14\,567 unités, 42~\% capturées, $+0.03$ ticks, 591 messages~; adjudication 2\,567, 0~\%, $-1.12$, 330~; au prorata
10\,884, 47~\%, $-0.17$, 474~; crypto 7\,267, 56~\%, $-0.07$, 2\,857~; événement 8\,367, 6~\%, $-0.52$, 799.
\item Prorata (taille 500~: 40\,633 unités, 58~\% capturées, $-0.14$, 419 messages)~; crypto (rythme de 16 s~: 8\,333 unités, 191 messages, aucun rejet)~;
événement (seuil de 3 ticks~: 7\,900 unités, 238 messages, $-0.57$).
\item Jusqu'à huit secondes, la plateforme accepte la totalité du budget, soit 151 à 154 ordres~; le volume croît à mesure que le rythme les espace, culmine à huit secondes et chute
au-delà de seize secondes, où le budget n'est plus utilisé.
\item Le tableau des résultats du chapitre~: volume, part capturée, edge et messages sur chaque plateforme, avec les paramètres d'origine et les paramètres propres.
\item Un paramètre~: prorata, crypto, événement. Rien~: transparente, inversée (la seule ligne de frais), à \omterm{def:hf:latency-arbitrage-and-its-defence:bump}{ralentisseur}, contrats à terme FIFO. Code différent~: adjudication et événement.
\item Le retard s'applique à l'annulation comme à l'ordre de l'arbitragiste~; un retard asymétrique qui épargnerait les annulations s'en écarterait.
\item Coter les nouveaux ordres en fonction du carnet tel qu'il sera à leur arrivée, ou n'en envoyer aucun tant que le cours bouge~; conserver les ordres au repos, que le retard
protège.
\item Coter autour d'une fair value plutôt qu'au meilleur niveau, et dimensionner selon le déséquilibre de l'adjudication.
\item Rédiger la moitié plateforme à partir du règlement de marché et des spécifications~; valider avec la moitié cotation d'origine~; faire tourner le simulateur avec les paramètres d'origine et
proposés~; l'environnement de test de la plateforme et les tests de conformité~; petite taille sous contrôles de risque~; pleine taille une fois que les exécutions, les rejets et les frais correspondent au
profil.
\item Le versionner~; chaque avis de plateforme donne lieu à une nouvelle version, validée et rejouée avant sa date d'effet~; le responsable du desk et le propriétaire de la plateforme donnent leur aval.
\item Apprendre les règles d'une plateforme avant le premier ordre plutôt que par les pertes.
\end{enumerate}
\end{solution}

\begin{solution}{iq:hf:a-venue-playbook:1}
Règle d'appariement, tick et lot, frais et unité, flux, types d'ordres, retards, budgets, séance et particularités, ainsi que les paramètres du desk. Avant le premier
ordre~: le lot (une taille que la plateforme refuse), le budget et ce que fait un dépassement, et le signe des frais.
\iqlookfor{règles liées aux paramètres, et non une liste de faits.}
\end{solution}

\begin{solution}{iq:hf:a-venue-playbook:2}
La taille~: l'allocation suit la taille affichée, si bien que la taille d'origine n'obtient qu'une petite part. Le risque est qu'une taille plus grande soit aussi davantage consommée lorsque
le cours saute (58~\% contre 47~\% du volume dans les simulations du chapitre), et que chacun surcote à son tour.
\iqlookfor{la taille comme paramètre du prorata, et son coût.}
\end{solution}

\begin{solution}{iq:hf:a-venue-playbook:3}
Le budget d'ordres~: les recotations à chaque mise à jour du carnet consomment les jetons par rafales, chaque ordre rejeté est renvoyé, et les jetons sont épuisés quand le marché
bouge. Rythmer les recotations dans le budget, laisser les annulations hors de lui, et dépenser la rafale sur les mouvements, non sur le bruit.
\iqlookfor{lire les rejets comme un problème de budget, et non comme une panne réseau.}
\end{solution}

\begin{solution}{iq:hf:a-venue-playbook:4}
Les cotations au repos peuvent être annulées avant l'arrivée du preneur de liquidité retardé, si bien que la sélection adverse subie par le fournisseur de liquidité face aux preneurs rapides diminue et qu'il peut coter plus serré
ou plus gros~; les preneurs rapides perdent, et, en partie, les preneurs lents qui attendent plus longtemps. Un fournisseur de liquidité dont les propres nouveaux ordres sont retardés paie en placements périmés.
\iqlookfor{qui paie pour la protection.}
\end{solution}

\begin{solution}{iq:hf:a-venue-playbook:5}
Au plus $b$ ordres par seconde en moyenne, si bien qu'une recotation de deux ordres toutes les $2/b$ secondes ne dépasse jamais le budget. Un rythme plus rapide peut générer davantage de transactions, parce que la
rafale $B$ absorbe les brèves poussées et les cotations restent plus près du meilleur niveau, au prix de rejets occasionnels.
\iqlookfor{débit moyen contre rafale.}
\end{solution}

\begin{solution}{iq:hf:a-venue-playbook:6}
Le meilleur cours acheteur du fournisseur de liquidité ne s'exécute que si le prix de règlement se situe à ce cours ou en dessous, ce qui arrive lorsque les ordres de vente de l'adjudication dépassent les ordres d'achat à
ce prix~: ses exécutions sont conditionnées à une pression vendeuse, si bien qu'elles sont suivies de cours plus bas. Une règle qui cote autour d'une fair value ajustée du
déséquilibre attendu de l'adjudication, et qui réduit ses tailles lorsque le déséquilibre est informatif, ne subit pas la même sélection adverse.
\iqlookfor{conditionnement sur l'événement de règlement.}
\end{solution}
